%======================================================================================
\section{Fission and fusion: where the sizes come from}\label{sec:fission}

$T(n)$ is the share of the time spent in configurations of $n$ pilots: the chance that a group seen at
a random instant has $n$ pilots (Figure~\ref{fig:describe}, left). If sizes change one pilot at a
time, the flux from $n$ to $n+1$ pilots balances the flux back, a birth--death chain
(Appendix~\ref{app:models}):
\begin{equation}\label{eq:chain}
  \frac{T(n+1)}{T(n)}=\frac{g(n)}{l(n+1)},
\end{equation}
with $g(n)$ the joins and $l(n)$ the leaves and landings per hour at size $n$ (Table~\ref{tab:rates}).
Splits and merges are left out. Per member, $g(n)=a\,n$ and $l(n)=\mu\,n^{1-\gamma}$: every member brings
in others at rate $a$ and leaves at rate $\mu n^{-\gamma}$, fitted to the rates, not to $T$. With
$\gamma=0$ the chain gives the log-series $T(n)\propto x^n/n$, $x=a/\mu$. For comparison, the form of
\cite{bonabeau1999}, $T(n)\propto n^{-\alpha}e^{-n/n_c}$, fitted to $T$.

\begin{figure}[H]
\centering
\includegraphics[width=0.9\textwidth]{fission.pdf}
\caption{Left: joins and losses per member per hour by size (mean size of each class; 8+ pooled),
with $a$ and $\mu n^{-\gamma}$ dashed. Right: $T(n)$ observed (sizes with at least 10~h) and from
the chain~\eqref{eq:chain} with the measured rates, with the per-capita rates and with $\gamma=0$;
dashed purple: the form of Bonabeau et al.\ fitted to $T$. Bands: 5--95\%.}
\label{fig:fission}
\end{figure}

\begin{reading}
\textbf{Reading.}
\begin{itemize}
  \item \textbf{One pilot at a time is enough.} With the measured rates the chain gives $T(n)$ within 4\%
        from 2 to 8 pilots (\StChainThreePred{} against \StChainThreeObs{} at 3 pilots, just outside its band),
        although splits and merges are \StMultiTwo\% of the changes at 2 pilots and \StMultiMax\% at 8+.
  \item \textbf{Joining grows with the members, leaving does not.} Each member brings in others at
        $a=\StCapitaA$ \StCapitaABand{} per hour, at every size. Each leaves at \StLossTwo{} per hour in a pair and
        \StLossMax{} at 8+: $\gamma=\StGamma$ \StGammaBand.
  \item \textbf{The heavy tail comes from staying.} With $\gamma=0$ ($x=\StX$) $T(\StNTop)$ is \StLogSeriesFactor{}
        times too small; with $\gamma=\StGamma$, \StCapitaFactor{} times. Fitted directly,
        $\alpha=\StAlpha$ and $n_c=\StNc$.
\end{itemize}
\end{reading}
